\section{Background and Related Work}
\label{sec:related}

\subsection{Disaggregated LLM serving and KV movement}
Prefill/decode disaggregation has become a standard abstraction: compute-intensive prefill and
memory-bound decode are placed on separate workers, and KV cache is moved between them over the
cluster network \cite{distserve,splitwise,mooncake}. Production systems push this further with
cache-centric global schedulers \cite{mooncake}, dynamic rescheduling and migration
\cite{llumnix}, phase-aware placement \cite{splitwise}, chunked prefill to smooth the
throughput--latency tradeoff \cite{sarathi}, and locality-aware loading \cite{serverlessllm}.
Because KV cache is large and its transfer sits on the critical path, the network side of this
architecture has become its own research area: KV compression and streaming \cite{cachegen},
compression and routing-collision avoidance, and network-aware placement of the two phases
\cite{helix,flowkv}. Recent position papers argue that the scarce resource has shifted from raw
FLOPs toward \emph{managed state, placement, network movement, and decision quality}
\cite{cloudabstractions,controlplane}. This paper takes that shift as given and studies the
\emph{decision quality} part: what the router must know, how fresh that knowledge must be, and
when acting on it is worse than ignoring it.

\subsection{Network-aware scheduling}
Network-aware scheduling has a published history, mostly outside LLM serving. CASSINI places
training jobs using periodically refreshed communication-pattern measurements \cite{cassini};
VClos isolates tenants to avoid ECMP hash collisions in multi-tenant GPU clusters \cite{vclos};
TopoOpt co-optimizes topology and parallelization \cite{topoopt}. For LLM serving specifically,
Helix solves a max-flow placement problem over heterogeneous GPUs and links \cite{helix}, and
GenSched jointly selects servers and allocates resources for prefill/decode stages in
heterogeneous edge networks \cite{gensched}. Concurrent work closest to ours is NetKV, which
selects decode instances using a network-cost oracle over a modeled dynamic congestion field
\cite{netkv}, and Aegis, which places work on an accelerator cluster using per-second switch
queue and ECN telemetry \cite{aegis}. Two observations motivate our study. First, all of these
works evaluate a specific mechanism in a specific configuration; none characterizes the
\emph{boundary} at which the network term stops carrying information. Second, none of them treats
the correlation between the compute ranking and the bandwidth ranking as an experimental
variable---Helix, for instance, solves for a placement directly rather than studying how much the
answer depends on the two rankings agreeing.

\subsection{Stale and delayed state information}
That stale state degrades load balancing is old knowledge: stale load data leads to wrong server
selection \cite{staleloadold}, delayed-information load balancing was formally listed as an open
problem \cite{delayedopen}, and stale link state causes routing to oscillate
\cite{stalelinkqos}. ParSync turned admissible staleness into a design variable by relaxing
freshness to gain scalability \cite{parsync}. The Age-of-Information literature, however, warns
against treating age itself as the objective: no universal staleness threshold exists, and the
mapping from age to harm is application-specific \cite{aoisurvey,timelyremote}. Telemetry-side
work optimizes freshness directly \cite{freshint}, and recent systems report that
freshness-based guards can be far too conservative---blocking $92$--$95\%$ of harmless cases in a
different domain \cite{staleguard}.

Within LLM serving, NetKV derives a bound under which a stale estimate does not change the
ranking, and reports a refresh-interval sweep whose conclusion is that the benefit is
insensitive \cite{netkv}; Aegis reports that gains hold under telemetry delays up to
\SI{60}{\second} \cite{aegis}. Our results are consistent with both, and explain them: both
sweeps vary only the link-side observable, which our measurements show is the slow-expiring
channel (\S\ref{sec:rq2}). We additionally (i) vary the queue-side observable, which expires
roughly an order of magnitude faster, (ii) sweep the two jointly rather than one at a time, and
(iii) separate \emph{latency} from \emph{time-averaging}, which are different distortions with
different consequences. Finally, our own first attempt at adapting to staleness---discounting
\emph{within} the link term---was falsified by our experiments (\S\ref{sec:discount}); we report
it because the failure mode (muting the congestion penalty makes ``fat link, slow compute''
workers look attractive) is easy to reproduce.

\subsection{Positioning}
Before positioning, we note the serving and fabric context our measurements sit inside. The
engines that disaggregated systems build on---PagedAttention/vLLM \cite{vllm}, SGLang
\cite{sglang}, Orca \cite{orca}---and the multiplexing strategies layered above them
\cite{alpaserve} define what ``a worker's queue'' means; industrial frameworks such as Dynamo
\cite{dynamo} and llm-d \cite{llmd} expose exactly the placement decisions we study; recent
systems push on KV locality and imprecise request metadata \cite{jitserve}, stale-state
handling in packet-level load balancing \cite{psp}, cross-region balancing \cite{gorgo}, and
freshness of RAG-side data \cite{churnbench}. The fabric itself is the standard one---Clos
topologies \cite{fattree} with congestion control \cite{dctcp,dcqcn,hpcc}---and published
inference throughput ceilings \cite{mlperf} set the scale at which queueing matters. Our study
does not modify any of these; it measures when a scheduler should include a network term at all.

Table~\ref{tab:positioning} summarizes the gap. Relative to prior work we (a) introduce a
configuration variable that predicts the magnitude of the benefit, (b) characterize staleness as
two channels with different expiry plus an averaging dimension, (c) evaluate \emph{conditional}
use of network information against both of its extremes and report where it loses, and (d)
measure the cost of \emph{obtaining} state on real containers, which simulation omits. Honest
negative results and boundary studies now have precedent in adjacent areas---self-described
negative results in KV caching \cite{elastickv}, design-space studies asking when each
disaggregation level pays off \cite{howfar}, equilibria that switch regime at saturation
\cite{poa}; we follow that template rather than the ``system plus headline speedup'' one.

Two scope clarifications close the positioning. First, what the reference scheduler
\emph{deliberately omits}: continuous batching, prefix-cache locality, KV-aware routing, and
prefill--decode phase scheduling are orthogonal mechanisms that real engines already implement
\cite{vllm,orca,alpaserve}, and shared-fabric effects (oversubscribed uplinks, switch queues,
incast) and concurrency-dependent service rates lie outside the model; the exact boundaries are
stated in \S\ref{sec:setup} and \S\ref{sec:threats}. We omit them because the question here is
whether a scheduler should include a network term \emph{at all}, and bundling them in would
entangle that question with mechanisms this study does not measure. Second, why
\textsc{Random}, \textsc{LeastQueue}, and \textsc{LeastLoad} suffice as base comparisons: they
span the three information regimes a scheduler can occupy---no state, queue state only, and
load state only---so any network term is evaluated as a strictly nested increment over the
third; our strengthened baselines (\S\ref{sec:rq1}) additionally compare against
power-of-two-choices, congestion-only, and bottleneck-minimizing scoring styles drawn from
recent practice. We characterize a boundary; we do not propose a scheduler that dominates
existing ones.

\begin{table}[t]
\centering
\caption{Positioning against the closest prior work. ``Channels'' = how many observed quantities
are swept for freshness; ``Boundary'' = whether the paper characterizes where the benefit
disappears or reverses; ``Align.'' = whether compute--bandwidth rank agreement is an experimental
variable.}
\label{tab:positioning}
\footnotesize
\setlength{\tabcolsep}{4pt}
\begin{tabular}{@{}lccccc@{}}
\toprule
Work & Venue & Serving & Chan. & Bound. & Align. \\
\midrule
CASSINI \cite{cassini}          & NSDI'24     & no  & 1 & no     & no \\
VClos \cite{vclos}              & ComNet'25   & no & 1 & no  & no \\
Helix \cite{helix}              & ASPLOS'25   & yes & 1 & no     & no \\
GenSched \cite{gensched}        & INFOCOM'26  & yes & 1 & no     & no \\
NetKV \cite{netkv}              & preprint    & yes & 1 & partial & no \\
Aegis \cite{aegis}              & SIGCOMM'26  & no  & 1 & no     & no \\
\textbf{this paper}             & ---         & yes & \textbf{2+win.} & \textbf{yes} & \textbf{yes} \\
\bottomrule
\end{tabular}
\end{table}
